\begin{table}
\centering
\caption{\label{tab:power_next_election}Precision of Next-Election Estimates}
\centering
\resizebox{\ifdim\width>\linewidth\linewidth\else\width\fi}{!}{
\fontsize{8}{10}\selectfont
\begin{threeparttable}
\begin{tabular}[t]{lllllll}
\toprule
State & Elections & N & Estimate (pp) & Cluster SE (pp) & 95\% CI (pp) & Approx. MDE (pp)\\
\midrule
Rajasthan & 2005--2010 & 4,061 & 1.4 & 0.9 & {}[-0.4, 3.3] & 2.6\\
Rajasthan & 2010--2015 & 3,867 & -2.8 & 1.1 & {}[-5.1, -0.6] & 3.2\\
Rajasthan & 2015--2020 & 4,014 & 0.4 & 1.1 & {}[-1.8, 2.7] & 3.2\\
Uttar Pradesh & 2005--2010 & 28,002 & 2.1 & 0.4 & {}[1.2, 2.9] & 1.2\\
Uttar Pradesh & 2010--2015 & 25,963 & 1.6 & 0.5 & {}[0.7, 2.6] & 1.3\\
\addlinespace
Uttar Pradesh & 2015--2021 & 30,975 & 3.6 & 0.5 & {}[2.6, 4.5] & 1.3\\
\bottomrule
\end{tabular}
\begin{tablenotes}
\item \textit{Note: }
\item Normal-approximation MDEs use 80\% power and two-sided 5\% size, computed from each fitted model or full-covariance linear contrast. They are precision summaries, not evidence of a zero effect or exact bootstrap power.
\end{tablenotes}
\end{threeparttable}}
\end{table}
